\documentclass[exercice,a4]{thedocument}%
\usepackage[french]{babel}%
\usepackage[np]{numprint}%
\usepackage{graphicx}%
\usepackage{enumitem}%
\startdatakeys
    \source{22GENMATMEAG1}
    \BO{2026}
    \cycle{4}
    \competencesDuSocle{CA,RA}
    \objectifsApprentissage{C4B3ca,C4C2aa}
    \motsCles{droites parallèles,Thalès,vitesse,durée,conversion}
\enddatakeys
\begin{document}%

Une famille se promène au bord d’une rivière.

Les enfants aimeraient connaître la largeur de la rivière.

Ils prennent des repères, comptent leurs pas et dessinent le schéma ci-dessous sur lequel les points C, E et D, de même que A, E et B sont alignés. (Le schéma n'est pas à l’échelle.)

\begin{center}
\eps[scale=0.5]{riviere-2}
\end{center}

\begin{enumerate}
\item Démontrer que les droites (AC) et (BD) sont parallèles.
\item Déterminer, en nombre de pas, la largeur AC de la rivière.
\end{enumerate}

 Pour les questions qui suivent, on assimile la longueur d'un pas à 65~cm.

\medskip

\begin{enumerate}[resume]
\item Montrer que la longueur CE vaut $13,3$~m, en arrondissant au décimètre près.
\item
	\begin{enumerate}
		\item  L’un des enfants lâche un bâton dans la rivière au niveau du point E. Avec le courant, le bâton se déplace en ligne droite en 5 secondes jusqu'au point C.

Calculer la vitesse du bâton en m/s.
	\item Est-il vrai que \og le bâton se déplace à une vitesse moyenne inférieure à \np[km/h]{10}\fg?
	\end{enumerate}
\end{enumerate}

\bigskip

\showthesource
\end{document}
